\subsection{Measured cautionary evidence from the structural-learning lane}
\label{sec:structural-cautionary}

The programme's structural-learning lane (its Paper IV artifacts, which remain intact and byte-unchanged in the research record) reached a terminal its own routing gate classified as conceptual cross-paper evidence rather than a standalone result. We absorb it here in that capacity, because both of its findings are direct measurements about the risks of \emph{method evolution itself} --- the subject of this paper. Nothing in this subsection promotes a structural-learning law; the defensible objects are two negatives and one instrument mechanic.

\paragraph{An adaptive method-evolution candidate lost to its static parent, with attributed cause.}
A learner-state-conditioned adaptive training allocator --- exactly the kind of candidate a method-evolution loop generates --- was compared against a static structural parent at strictly equal example budget across six frozen simulation worlds (384 replicates). The adaptive arm lost on balanced mastery by $-0.0166$ (bootstrap 95\% CI $[-0.0174, -0.0158]$) and degraded the hard-safety floor. Failure attribution then localized the loss to the allocation-policy stage via two designed lever arms: \emph{guard-rail budget capture} (a coverage floor implemented as a budget-consuming target rather than a constraint absorbed 13 of 48 examples into the highest-mastery coordinate in every world; capping it recovered about a third of the gap) and \emph{within-round concentration} (near-uniform budget still lost $-0.0108$). One diagnostic arm was found confounded and is recorded as uninformative rather than used as evidence. The originally recorded root-cause narrative of the parent receipt was refuted by the realized budget counts and is superseded \emph{as an interpretation only}; the receipt itself is immutable and preserved. Both identified defects were verified present in the promoted production allocator, so the diagnosis transfers beyond the simulator. The operational consequence stands: the static parent retains default authority, and a preserved negative about our own adaptive mechanics is part of this paper's method-evolution evidence, not an embarrassment to be routed around.

\paragraph{The instrument itself could not have adjudicated the question --- provably.}
The deeper finding is about the evaluation instrument. A three-tier ceiling analysis showed that \emph{no} equal-budget allocation policy, however optimal, could have reached the instrument's own registered material-effect gate of $0.05$: a greedy-oracle policy achieved $+0.0015$ (a lower bound), local search found $+0.0045$ (a lower bound), and a rigorous harm-free-relaxation upper bound caps every policy at ${\approx}0.0246$ --- a factor of two below the gate, stated with its tightness caveat because an earlier, stronger version of this bound was itself found invalid and corrected before use. Conventional power analysis was silent throughout: bootstrap intervals were ${\sim}0.0016$ wide, amply ``powered'' for an effect the instrument could not generate. An independent lane that tuned allocation policies on the same instrument class plateaued at the constructive ceiling computed here and stopped --- corroborative, cross-instrument agreement between a mechanism bound and an empirical search exhaustion.

\paragraph{What this changes for governed method evolution.}
These are two instances of one lesson, and it is the same lesson taught by the gate-falsifiability audit of Section~\ref{sec:evaluation}: before a confirmatory budget is spent, the evidence instrument must itself be qualified. Falsifiability audit asks whether a gate \emph{can fail}; admissibility analysis asks whether the instrument \emph{can express the effect it is gating}. The structural-learning lane contributes the second mechanic in executable form: compute the instrument's oracle ceiling before an allocation comparison and require it to clear a pre-frozen multiple of the minimum detectable effect, emitting \path{INSTRUMENT_INADMISSIBLE_CEILING_BELOW_GATE} --- and spending nothing --- otherwise, failing closed to \CANNOTCHECK{} when the oracle is not computable. This paper's governed upgrade protocol adopts that requirement prospectively for its own evaluation epochs: a method-evolution challenger may not consume a confirmatory epoch on an instrument whose ceiling has not been shown to clear the registered effect gate. The fresh-task lift protocol of this paper is explicitly flagged by the same logic, and applying the mechanic to it returned the fail-closed branch rather than a bound (\path{research/paper3_lift_ceiling_qualification_v1/}). Two things had to be corrected in the flag itself. First, the $0.05$ carried in that protocol is a two-sided significance level---it appears there only as $\alpha$ and as the Holm family-wise rate, beside a planning effect of roughly $15$ percentage points---so it was never an effect gate, and comparing an oracle ceiling against it compares a magnitude with a probability. The lift lane has \emph{no} registered material-effect threshold, which is a stronger statement than having one that is unqualified. Second, the ceiling is not merely unrecorded but not derivable before execution for this instrument class: the oracle requires a generative model of per-arm outcomes for the pinned subject, obtainable only from pilot data on the very run being gated, and the frozen evaluator hash and task panel are likewise unset. The terminal is \path{CANNOT_CHECK__ORACLE_NOT_COMPUTABLE_NO_GENERATIVE_MODEL}, invariant across an $18$-cell sweep of candidate effect thresholds and clearance multiples, with the sibling lane's published bound reproduced to $1.5\times10^{-16}$ as a positive control and a second control returning \path{ADMISSIBLE}, so the outcome is a measurement rather than a default.

The structural reason is worth stating, because it bounds where the mechanic applies. The sibling lane's four allocators were \emph{policies} over a shared example budget, so its water-filling relaxation had an allocation space to relax; this paper's four arms are fixed \emph{conditions}, and no such space exists. An oracle-ceiling qualification is therefore available for allocation instruments and unavailable for condition-contrast instruments, whatever their budget. The practical consequence for the blocked obligation is narrow and unwelcome: the qualification cannot de-risk the accelerator spend in either direction, so the block is not reducible to budget. What is local, free and prerequisite is freezing the evaluator hash and a material-effect threshold for the lift; until those exist there is no gate for a ceiling to clear.
